\chapter{Marco Teórico}

El desarrollo de soluciones tecnológicas para la gestión sostenible de residuos orgánicos exige una comprensión integrada de los procesos biológicos, ambientales y computacionales involucrados. Este capítulo ofrece la base teórica que sustenta la propuesta de un sistema inteligente híbrido para la predicción de emisiones de gases de efecto invernadero (GEI) en la bioconversión con BSF. Se articulan dos ejes interdependientes: (i) la bioconversión como estrategia de valorización de residuos en contextos agroindustriales; (ii) los enfoques de modelado matemático y sistemas inteligentes para la estimación dinámica de emisiones.

\section{Residuos orgánicos y bioconversión con \textit{Hermetia illucens}}

En contextos agroindustriales del sur global —como el Valle del Cauca—, los residuos orgánicos provenientes de cultivos como café, caña y frutas se generan de forma dispersa, en picos estacionales y con alta carga biodegradable, lo que incrementa el riesgo de emisiones de CH\textsubscript{4} cuando su gestión es inadecuada \parencite{FAO2021,Leddin2024}. Frente a las limitaciones de las estrategias convencionales —rellenos sanitarios, digestión anaerobia y compostaje—, que suelen ser inviables técnica o económicamente en estos entornos \parencite{kongEvaluatingGreenhouseGas2012,nordahlGreenhouseGasAir2023}, la bioconversión con BSF se ha adoptado crecientemente como una alternativa de economía circular.

Esta tecnología transforma residuos en biomasa larval rica en proteínas y lípidos, y en \textit{frass}, un fertilizante orgánico de alta calidad, generando valor en lugar de desechos \parencite{Surendra2016,bermudezComprehensiveUtilizationBlack2023,singhInclusiveApproachOrganic2019}. Aunque estudios preliminares sugieren que reduce emisiones de CH\textsubscript{4} y N\textsubscript{2}O frente al compostaje \parencite{jenkinsProcessingPoultryManure2025}, su implementación carece de sistemas rigurosos para cuantificar y predecir las emisiones asociadas en tiempo real, lo que limita su validación como solución climática verificable.


Durante la bioconversión de residuos orgánicos, la generación de dióxido de carbono (CO\textsubscript{2}) proviene principalmente de la respiración metabólica tanto de las larvas  como de la microbiota asociada al sustrato. Las larvas ingieren materia orgánica y, mediante procesos catabólicos aeróbicos, oxidan los compuestos carbonados para obtener energía, liberando CO\textsubscript{2} como subproducto principal. Estudios indican que hasta un tercio del carbono ingerido puede ser emitido como CO\textsubscript{2} \parencite{dienerConversionOrganicMaterial2009}. Adicionalmente, la actividad microbiana en el sustrato —especialmente durante las fases iniciales de degradación— contribuye significativamente a las emisiones gaseosas. Factores como la composición del sustrato, la densidad larval, la temperatura, la humedad y la oxigenación influyen directamente en la tasa de respiración y, por tanto, en las emisiones totales de CO\textsubscript{2} \parencite{rossiEstimatingDynamicsGreenhouse2024,bekkerImpactSubstrateMoisture2021a,eriksenDynamicModellingFeed2022}. A diferencia del metano (CH\textsubscript{4}), el CO\textsubscript{2} generado en este proceso proviene mayoritariamente de rutas aeróbicas, lo que lo convierte en un indicador sensible de la actividad biológica en tiempo real.

La aptitud de los residuos para la bioconversión varía significativamente según su composición. Como se resume en la Tabla~\ref{tab:residuos_BSF}, subproductos agroindustriales como la pulpa de café o cáscaras de frutas son altamente adecuados, mientras que residuos lignocelulósicos o lodos orgánicos requieren pretratamiento o presentan riesgos sanitarios \parencite{elsayedConversionProteinrichWaste2024,brunoValorizationOrganicWaste2025}. Esta heterogeneidad —común en el Valle del Cauca— subraya la necesidad de sistemas de monitoreo adaptables, capaces de caracterizar el proceso en función del sustrato y las condiciones locales.

La evidencia cuantitativa disponible es fragmentada: se reportan emisiones de CO\textsubscript{2} entre 7.76 y 11.88\,g por gramo de larva seca \parencite{rossiEstimatingDynamicsGreenhouse2024}, o 96.5\,g de CO\textsubscript{2} por kg de residuo tratado \parencite{fuhrmannComprehensiveIndustryrelevantBlack2025}, pero sin protocolos estandarizados ni mediciones continuas. Esta incertidumbre impide comparar tecnologías, diseñar políticas públicas informadas o acceder a mecanismos de mercado basados en la mitigación climática.

\begin{table}[H]
\centering
\caption{Clasificación de residuos orgánicos y su aptitud para bioconversión con \textit{Hermetia illucens}}
\label{tab:residuos_BSF}
\begin{tabular}{|p{4cm}|p{5cm}|p{5cm}|}
\hline
\textbf{Tipo de residuo} & \textbf{Origen común} & \textbf{Aptitud para \textit{BSF}} \\ \hline
Residuos agrícolas & Restos de cosecha, podas, hojas secas & Moderada. Requieren pretratamiento por bajo contenido proteico. \\ \hline
Residuos alimentarios & Sobras de cocina, frutas, verduras, pan & Alta. Contenido balanceado de lípidos y carbohidratos. \\ \hline
Subproductos agroindustriales & Pulpa de café, cáscaras de frutas, bagazo de caña & Alta. Composición variable, pero generalmente adecuada. \\ \hline
Grasas y aceites residuales & Aceites usados, residuos grasos de cocina & Media. Requieren mezcla con otros residuos para evitar inhibición. \\ \hline
Lodos orgánicos & Lodos de depuradoras, biosólidos & Baja. Riesgo sanitario y composición inconsistente. \\ \hline
\end{tabular}
\end{table}

\textbf{Hacia sistemas de monitoreo adaptable en contextos heterogéneos:}

La heterogeneidad de residuos agroindustriales documentada en la Tabla~\ref{tab:residuos_BSF} 
—desde pulpa de café (alta aptitud) hasta lodos orgánicos (baja aptitud)— refleja una realidad 
común en el Valle del Cauca: los productores raramente procesan un único tipo de residuo de forma 
consistente. En su lugar, enfrentan mezclas variables según la estacionalidad, disponibilidad local 
y dinámicas de producción \parencite{elsayedConversionProteinrichWaste2024,brunoValorizationOrganicWaste2025}. 
Adicionalmente, variables ambientales críticas —temperatura, humedad del sustrato, oxigenación— 
fluctúan durante el ciclo de bioconversión, particularmente en contextos rurales donde el control 
es limitado \parencite{bekkerImpactSubstrateMoisture2021a,rossiEstimatingDynamicsGreenhouse2024}. 

Esta variabilidad operativa implica que los métodos de medición estática o puntual —como 
muestreos puntuales con cromatografía de gases— son insuficientes para caracterizar emisiones 
de manera confiable. Se requieren sistemas de monitoreo que sean: (i) \textbf{continuos en tiempo 
real}, capaces de capturar dinámicas horarias y diarias; (ii) \textbf{adaptables a diferentes tipos 
de residuos}, disminuyendo la necesidad de recalibración manual frecuentemente; (iii) \textbf{interpretables}, 
permitiendo a productores locales comprender qué sucede en el sistema y tomar decisiones operativas 
informadas. Es precisamente esta necesidad la que motiva un enfoque integrado que combine 
modelación matemática dinámica, sensores de bajo costo y calibración automática mediante algoritmos 
de aprendizaje automático, tema de la siguiente sección.

\section{Modelación de emisiones y sistemas inteligentes}

La cuantificación dinámica de emisiones en procesos biológicos requiere 
herramientas que combinen comprensión mecanística y adaptabilidad empírica. 
En este sentido, el modelado matemático basado en ecuaciones diferenciales 
ordinarias (EDO) ha facilitado avances significativos. Por ejemplo, modelos 
de crecimiento larval como el de Eriksen (2022) integran balances energéticos 
y cinética logística para simular biomasa y producción de CO\textsubscript{2}, 
combinando principios del \textit{Dynamic Energy Budget} con eventos clave del 
ciclo de vida larval (muda y metamorfosis) \parencite{eriksenDynamicModellingFeed2022}. 
Similarmente, funciones cinéticas tipo Monod y Haldane —adaptadas de modelos 
de digestión anaerobia— han permitido describir degradación de sustrato y 
crecimiento microbiano en sistemas con insectos \parencite{winAnaerobicDigestionBlack2018,
ahlamineMathematicalAnalysisAnaerobic2024}. 

Sin embargo, estos modelos presentan limitaciones críticas cuando se implementan 
en contextos operativos reales: (i) la mayoría fue validada en condiciones 
de laboratorio controladas, sin integración de datos en tiempo real; (ii) generalmente 
no capturan el efecto simultáneo y dinámico de múltiples variables ambientales 
(temperatura fluctuante, humedad variable, oxigenación heterogénea) sobre las 
emisiones de CO\textsubscript{2} y CH\textsubscript{4}; (iii) dependen de parámetros 
que requieren calibración manual y periódica, sin mecanismo de actualización automática 
\parencite{bekkerImpactSubstrateMoisture2021a,rossiEstimatingDynamicsGreenhouse2024}. 
Estas limitaciones impiden que los modelos actuales predigan con confianza el 
comportamiento del sistema bajo variaciones operativas típicas en contextos 
agroindustriales rurales.

Paralelamente, los avances en sensores de bajo costo y plataformas IoT han 
facilitado el monitoreo continuo de variables críticas —CO\textsubscript{2}, 
CH\textsubscript{4}, temperatura, humedad— en entornos operativos reales 
\parencite{rajRealTimeEstimation2023,duobieneDevelopmentWirelessSensor2022,
yavariArtEMonArtificialIntelligence2023}. Estos datos, procesados mediante 
algoritmos de aprendizaje automático (AAA), permiten identificar patrones 
complejos y no lineales que serían difíciles de capturar mediante reglas 
explícitas \parencite{weichertReviewMachineLearning2019,chenOptimizationModelProcess2021}. 

En particular, la integración de modelos mecanísticos (EDO) con AAA ha demostrado 
potencial en sistemas biotecnológicos complejos: mientras el modelo EDO proporciona 
una estructura interpretable y basada en principios físicos, el AAA permite calibración 
dinámica de parámetros a partir de datos en tiempo real, mejorando predicciones cuando 
las condiciones operativas se desvían del escenario de entrenamiento 
\parencite{kobelskiModelbasedProcessOptimization2024a,guillaumeAsymptoticEstimatedDigestibility2023}. 

Sin embargo, esta integración enfrenta desafíos prácticos: (i) los algoritmos de AAA 
requieren suficientes datos históricos para entrenar sin riesgo de \textit{overfitting}, 
pero ciclos de bioconversión típicos generan pocas muestras por evento; (ii) la validación 
cruzada y selección de hiperparámetros requieren cuidado especial en contextos con datos 
limitados; (iii) la elección entre arquitecturas simples (regresión lineal múltiple) versus 
complejas (redes neuronales) dependerá del volumen y calidad de datos disponibles. Por ello, 
esta tesis adopta un enfoque escalonado: comenzar con modelos interpretables (regresión múltiple), 
evaluar redes neuronales multicapa solo si se dispone de base de datos amplia, y validar todas 
las predicciones contra mediciones de referencia de laboratorio.

En síntesis, el marco teórico desarrollado en este capítulo fundamenta un enfoque 
técnico híbrido que integra tres componentes interdependientes: (i) un modelo mecanístico 
basado en ecuaciones diferenciales ordinarias que captura la dinámica biológica del proceso 
de bioconversión; (ii) sensores IoT de bajo costo que permiten monitoreo continuo de emisiones 
de CO\textsubscript{2} y CH\textsubscript{4} en tiempo real; (iii) algoritmos de aprendizaje 
automático que facilitan calibración dinámica y predicción adaptativa bajo variaciones 
operativas. Esta integración es particularmente relevante para contextos agroindustriales 
del sur global, donde la falta de infraestructura costosa y la necesidad de soluciones 
escalables hacen que un sistema de bajo costo y mantenible sea más viable que tecnologías 
convencionales.

Operacionalmente, la ``inteligencia'' del sistema reside en su capacidad para: (1) generar 
predicciones con fundamento mecanístico (interpretables, no cajas negras); (2) actualizar 
predicciones automáticamente a partir de datos en tiempo real; (3) proporcionar información 
clara a productores locales para que tomen decisiones informadas sobre condiciones operativas 
(temperatura, humedad, manejo de residuos). No obstante, esta investigación reconoce 
explícitamente que la mera disponibilidad de datos técnicos no garantiza equidad o sostenibilidad. 
La verdadera gobernanza de un sistema como este requiere —más allá del alcance de esta tesis— 
decisiones sobre acceso a datos, beneficios compartidos, y co-diseño con usuarios finales. 
Estas dimensiones ético-políticas constituyen un campo de investigación complementario 
igualmente importante.